\section{Bố cục của đồ án}

Bố cục của đồ án được trình bày thành 5 chương. Mỗi chương trình bày và thảo luận các vấn đề khác nhau liên quan đến đồ án và được tóm tắt như sau:

\textbf{Chương 1: Giới thiệu.} 

Tổng quan về các vấn đề liên quan đến đồ án, những thông tin về robot và hệ thống đa robot, tìm hiểu về bài toán xây dựng đội hình, lên chiến lược đẩy vật và các phương pháp ước lượng vị trí điểm trọng tâm, các nghiên liên quan và thảo luận.

% \textbf{Chương 2: Webots và E-puck robot hai bánh vi sai} 

% Trình bày tổng quan về trình mô phỏng vật lý Webots  và cấu hình chi tiết của mô hình robot E-puck. Trình bày mô hình điều khiển động học hai bánh vi sai.

\textbf{Chương 2:  Cơ sở lý thuyết} 
Giới thiệu về các cơ sở lý thuyết được lấy làm cơ sở cho việc xây dựng đội hình, lên chiến thuật đẩy và ước lượng trong tâm trong đồ án bao gồm lý thuyết về grasping, mô hình ma sát Coulomb, mô hình đông lực học, thuật toán PSO và biến thể PSO phân tán, phương pháp Lagrange và Lagrange tăng cường, bài toán QP.   

\textbf{Chương 3:  Phương pháp đề xuất} 
Phân tích và làm rõ yêu cầu của hệ thống, từ đó lý giải chi tiết về cách xây dựng hệ thống và phát triển, mô hình hóa và áp dụng các thuật toán để giải quyết yêu cầu của các bài toán con.

\textbf{Chương 4: Mô phỏng thực nghiệm và đánh giá.} 
Chương này đề cập đến cách thức thiết lập các thí nghiệm sau khi đã xây dựng được hệ thống mô phỏng hoàn chỉnh. Các kết quả mô phỏng của sẽ được đưa ra để đánh giá dựa trên các thước đo đã thiết kế từ trước.

\textbf{Chương 5: Kết luận và hướng phát triển.} 
Các tổng kết về các công việc đã thực hiện trong đồ án và các kết quả đạt được. Đồng thời cung cấp các hướng phát triển trong tương lai dựa trên đồ án.